\documentclass{ltjsbook}
\usepackage{docmute} 
\usepackage{../../myStyle} 

\begin{document}
% \chapter{モデル比較}
\chapter{Model Comparison}
\label{x30_comparison}
% さて，今回でこの授業も最終回となりました。今回は授業全体のまとめとして，これまでの流れを俯瞰的に捉え直すとともに，今後みなさんがこの講義を通じて学んだことが，どのように利用できるのかに言及していきたいと思います。
So, this class has come to its final session. This time, as a summary of the whole series, while revisiting the flow of what we've been learning, I would like to mention how you all can utilize what you have learned through this lecture in the future.

% 本書の前半，心理学データ解析応用1として前期にお話ししてきた内容は，心理学で利用されるデータ，とくに心理尺度を用いたものが，どういう根拠で「心」を測定していると言えるのかについて解説しててきました。
The first half of this book, as discussed in the first term in Applied Psychology Data Analysis 1, explains how the data used in psychology, especially those using psychological scales, can be said to measure the 'mind' based on what grounds.
% 後半，心理学データ解析応用2の方は，プログラミングやベイズ統計といった目新しい研究手法の方が目についたかもしれませんが，その背後にあったメッセージは，どういうメカニズムでデータが生成されてきたのかという，\textbf{データ生成メカニズム}というアイデアをもつことでした。その意味で，本書全体のメッセージは一貫しています。すなわち，心理学者がデータを取るという時の，データに込められた意味や想定されるメカニズムをどのように表現し抽出し解釈するかという，心理学的営みそのものをお伝えしたかったのです。ここに無知・無自覚なまま，心理尺度で何かが測定できると考えたり，検定の結果から意味を解釈したりできるはずがないのですから。
In the latter half, the section on Applied Psychological Data Analysis 2, you may have noticed new research methods such as programming and Bayesian statistics, but the underlying message was to have the idea of a **data generation mechanism**, which is how the data came to be generated. In that sense, the message of the entire book is consistent. That is, when psychologists collect data, we wanted to convey how they express, extract, and interpret the meaning and assumed mechanism embedded in the data, which is the very activity of psychology itself. Without understanding this, one cannot be expected to believe something can be measured by a psychological scale, or to interpret meaning from test results.

% そのための準備として，確率統計に関する理論や，コンピュータ，プログラミングに関する技能などの習得が必要だったわけです。心理学を学ぼうと思って入ったはずの大学で，どうして数学やプログラミングをしているのか，こんなはずじゃなかったと思った人もいるかもしれませんが，逆にどうして既知の知識や技術だけで，この摩訶不思議な心，人間，社会のありようが理解できると思っているのかと言いたいほどです。もちろん他のアプローチや，他にも学ぶべきことがたくさんありますし，時間は有限ですから苦手な手法は後回しにしたくなる気持ちはわかりますが，実は他の教養を身につけるよりもこちらの方が直接的で近道なルートだといえるかもしれません。
As preparation for that, it was necessary to acquire skills in theory relating to probability statistics, computers, and programming. Some people may wonder why they are doing math and programming at a university where they supposedly enrolled to study psychology and think this was not what they had in mind. On the contrary, one might even say, how can you understand the mysterious nature of the mind, human beings, and society with just the knowledge and techniques you already have? Of course, there are many other approaches and things to learn, and time is limited, so I understand the desire to put off methods you're not good at. However, it could be said that this is a more direct and shortcut route than acquiring other knowledge.

% \section{ベイジアンモデリング}
\section{Bayesian Modeling}
% Stanを使った統計的アプローチは，\keyterm{ベイジアンモデリング}{ベイジアンモデリング}{Bayesian Modeling}と呼ばれることがあります。Stanが事後分布からの乱数発生機であり，事前分布と尤度を記述するだけで事後分布を得ることができるというところが，ベイズの定理に根拠を持ったベイズ統計学ですから，「ベイジアン」という冠がつくのは当然です。しかし後半の「\keytermL{モデリング}{もでりんぐ}」については，なにもベイズ統計学に限った話ではありません。モデルを立てて考えるということだけを考えれば，広い意味では心理学だって「心」というモデルを立てて考えているのですから。もっというと，学問というのは一般的に抽象化(あるいは単純化)した理想の世界での論理的構造を扱うものですから，すべてモデリングアプローチだと言えるかもしれません。
Statistical approaches using Stan can sometimes be referred to as Bayesian Modeling. It's natural to label it 'Bayesian' because Stan is a random number generator from the posterior distribution, and simply by describing the prior distribution and likelihood, we can obtain the posterior distribution - this is based on Bayes' theorem of Bayesian statistics. However, the latter term 'modeling' is not limited to Bayesian statistics. If we only consider the act of formulating a model, even psychology, in a broad sense, is creating a model of the 'mind'. Moreover, as academic research generally deals with logical structures in an abstracted or simplified ideal world, it could be said that all disciplines follow a modeling approach.

% 心理学はその研究方法として，調査，実験，観察，介入といったものがあります。これらを通じて心とはなにかということを考えていくわけです。とくに実験や介入においては，条件の制御や統制を行って，それに伴って生じた結果の違いを\textbf{効果}と考えるのでした。そのことと\keytermL{帰無仮説検定}{きむかせつけんてい}は非常に相性が良かった事は，想像に難くないでしょう。帰無仮説検定は，出てきたデータと変数の関係について，機械的に判断を下すことができる方法論でした。帰無仮説検定が農学から生まれてきたことからもわかるように，その方法はどの研究分野に対しても適用できる，共通のルール・御作法なのです。心理学者は\underline{実験計画に心理学的要素を埋め込み}，結果を統計的に判断することで，心理学的な結果の解釈を進めるてきました。いいかえれば，実験の計画を立てることそのものが心理学のエッセンスであったのです。
Psychology employs various research methods such as survey, experiment, observation, and intervention. Through these, we contemplate what the mind is. Particularly in experiments and interventions, differences in results that occur with the control and regulation of conditions are considered the "effect". It's not hard to imagine that these principles pair exceptionally well with null hypothesis testing. Null hypothesis testing is a methodology that allows for a mechanical judgement regarding the relationship between the data that has surfaced and its variables. As you can see from the fact that null hypothesis testing was born out of agriculture, its methodology can be applied to any research field, a common rule and etiquette, so to speak. Psychologists have "embedded psychological elements in experimental design", used statistical judgement for the results, and proceeded with the interpretation of psychological results. In other words, the planning of the experiment itself was the essence of psychology.

% しかし皆さんはすでに，実験計画が統計的に見れば\keytermL{線形モデル}{せんけいもでる}にすぎないことを知っていますね。そしてごく単純な線形モデルによってしか，心理学のエッセンスを表現する方法がなかった時代にあっては，想定しうる心のメカニズムが平均値の差として出てくるように置き換えるほかなかったのです。方法論に自由がない分，その他のところで工夫するしかなかったとも言えます。そうした工夫の中には芸術的なまでに作り込まれているものがあったりしますから，そこに痺れる憧れる，という人も少なからずいるのは当然ともいえるでしょう。
But you already know that, statistically speaking, an experimental plan is nothing more than a linear model, right? And in an era where the essence of psychology could only be expressed through a very simple linear model, the supposed mechanism of the mind had no choice but to be converted into differences in mean values. In the absence of freedom in methodology, it could be said that the only option was to be inventive elsewhere. There are some who are naturally attracted to the tingling feeling that comes from the artistic intricacy within such inventions.

% 線形計画の結果は，操作/介入の効果があったかなかったかという判断になりがちです。結果がそれしか示しておらず，その結果を生むための工夫はすでになされており，得られるデータの精度も平均値差を示すことができれば良い程度にしかない，ともいえるかもしれません。しかし測定のツールは時事刻々と進歩してきましたので，より精緻に違いを見て，よりリッチな解釈を許す意味合い豊かなデータも得られるようになってきました。そのような時，伝統的な方法論だけで立ち向かうのはもったいないのではないでしょうか。
The results of linear programming tend to lead to judgments of whether or not there was an effect of operation/intervention. The results only indicate that, and the efforts to produce those results have already been made. The accuracy of the data obtained is only good enough if it can show average difference. However, measuring tools have been advancing moment by moment, so we are now able to obtain more nuanced data that allows for more rich interpretations. At such times, it might be a waste to confront them with only traditional methodologies.
% 方法論の呪縛から自由になり，さまざまな表現ができるようになれば，心理学ももっと発展させることができるでしょう。
By freeing ourselves from the constraints of methodology and enabling various expressions, we might be able to further develop psychology.

% \keytermL{モデリング}{もでりんぐ}は，変数間関係を数式で表現することでもあります。「心」や「気持ち」といった表現を始め，「幸せ」とか「恋愛感情」といった日常用語で考えるべきところを，$X_i,Y_{ij},\alpha,\beta...$といった記号で置き換えて表現するのには慣れが必要です。ひょっとしたら抵抗を感じる人もいるかもしれません。しかし，数式で表現する事は，他の解釈を許さない一意的表現をすることでもあります。人間って色々だよね，当てはまることもあれば当てはまらないこともあるよね，という日常的な理解で留まりたいのであれば--そしてそこで止まっていられるのは幸せなことでもあるのですが--それでいいのですが，学問を進める以上はより明確に，より誤解のない表現をしなければなりません。個別のケースにしか当てはまらないこと，その程度を確率で表現して，理論的に全体的な傾向を「こうなっているに違いない」と厳密に表現するには，日常用語ではなく数式用語が必要なのです。モデリングを進める上で，そうした表現ができるようにトレーニングする必要はありますが，使えるようになると折線回帰や状態空間モデルのように，さまざまな表現ができることも見てきた通りです。
Modeling can also be the expression of relationships between variables in mathematical formulas. It can require some adjustment to replace everyday terms such as "heart" or "feelings", and concepts like "happiness" or "romantic emotions" with symbols like $X_i,Y_{ij},\alpha,\beta...$. There may be people who resist this. However, expressing things in mathematical form also means conveying them unambiguously, without room for alternate interpretations.

If you prefer to remain within the casual understanding of, "Well, people are diverse, so some things fit while others don't" -- and staying in that realm can be blissful -- then that's fine. But to advance knowledge, you must articulate things more clearly, in a way that avoids misunderstanding. Things that only apply to specific cases, expressing the extent of that applicability in probabilities, to rigorously claim that the overall trend must be "just so" - requires mathematical language, not everyday terms.

In advancing your modeling, it's necessary to train yourself to use this language. Once you master it, as we've seen, you'll have myriad ways to articulate things, like with piecewise regression or state-space models.


% この授業を受けてきたからと言って，さあ今すぐ自分で心理学的モデルを作れと言われても難しい，と思うかもしれません。実際これまで私が教えてきた中で，よく質問されるのは「やり方はわかるけれども，自分でやれる気がしない。どうやってモデルを思いついたらいいのか」ということです。これに対する答え方は2つあって，1つは「色々なパターンを見て模倣する」です。「学ぶ」は「真似ぶ」からきているとも言われたりしますが，どのような表現の可能性があるのかについては，実際の応用パターンをさまざま見るのが早いでしょう。馴染みのない料理や調理法があれば，レシピ本を買って応用パターンをいろいろ仕入れますよね。そうすることで，「こんな工夫ができるのか。じゃあ自分はこうしてみよう，ここをちょっと変えてみよう」と前に進むことができるのですから。こうしたモデルがたくさん載っているものとして，\textcite{Lee201709}や\textcite{Toyoda2017,Tanobei2018,Tanobei2019}などがありますので，パラパラっとみながら応用例を広く眺めてみるといいでしょう。
You might think it would be difficult to immediately create your own psychological model just because you've been taking this class. In fact, a common question I often get asked in the classes I've been teaching is: "I understand how to do it, but I don't feel like I can do it myself. How can I come up with a model?". There are two ways to answer this; one is "see various patterns and imitate them". It is often said that "learning" comes from "imitating", and the fastest way to know what kind of expression possibilities there are is to look at a variety of practical application patterns. If you're unfamiliar with certain cuisines or cooking methods, you'd buy a recipe book and stock up on a variety of application patterns, right? By doing so, you can make progress by thinking, "So this kind of ingenuity is possible, let's try it this way, let's change this part a bit". There are resources like \textcite{Lee201709} or \textcite{Toyoda2017,Tanobei2018,Tanobei2019} which have many such models, so it would be good to casually browse through a wide range of application examples.

% もう1つの答えは，「紙とペンをつかって設計図を書く」というものです。この方法は，この講義の中で再三お伝えしてきたやり方になります。データを見た時に，どういうメカニズムでデータが得られたかを考えるところにこそ，心理学者の本当の興味関心があるはずです。このデータは心のこういう仕組みを反映しているのではないか。データの背後には変数同士のこういう関係が潜んでいるのではないか，ということを考えていくわけです。わからないところは確率で表現しつつ，変数間関係を関数で記述する。そのための設計図です。設計図とStanがあれば答え(事後分布)は得られるのですから，データ生成の仕組みを作り込むことに我々は全力を傾けることができます。複雑な関数関係をいきなり思いつくのが難しい，ということもあるでしょう。その場合はデータを可視化し，データに合うような関数はどんな形なのかを考える，あるいはいったん，単純な線形モデルを当てはめて，それを徐々に複雑にしていく，というやり方が良いでしょう。\textcite{Maeda2017}は心理学者が書いたベイズ統計の本ですが，その後半は線形モデルです。一般化線形モデルや階層モデルについて，とくに複雑な前置きなく「こういう仕組みになってるんだから，とりあえずそれを反映して，関係は線形でも大体うまくいくでしょ」という感じで話が進んでいきます。簡単なモデルから徐々に複雑にしていく，というのがモデリングの王道でもあります。完成品だけをみるととても複雑に思えますが，設計図片手に紐解いていけば，意外とわかりやすいかもしれません。
Another answer is to "draw a blueprint with paper and pen." This is the method that I have repeatedly mentioned in this lecture. When you see data, where a psychologist should really be interested is in thinking about what kind of mechanism the data was obtained. This data may reflect some kind of mechanism of the mind. It means thinking about whether there might be such a relationship between variables behind the data. While expressing the unknown parts with probability, describe the relationship between variables with a function. This is a blueprint for that. If you have a blueprint and Stan, you can get the answer (posterior distribution), so we can devote ourselves to developing the mechanism of data generation. It might be difficult to come up with complex function relationships all at once. In that case, visualizing the data and thinking about what shape a function that suits the data would have, or first applying a simple linear model and gradually making it more complex, is a good way to do it. \textcite{Maeda2017} is a book on Bayesian statistics written by a psychologist, and the second half is about linear models. It goes on to talk about generalized linear models and hierarchical models, especially without a complicated preface, saying "Since this is how it works, just reflect that, and the relationship will probably work well even if it is linear." It is also a royal road to modeling to gradually complicate from a simple model. If you only look at the finished product, it may seem very complicated, but if you unravel it with a blueprint in hand, it might be surprisingly easy to understand.

% \keytermL{モデリング}{もでりんぐ}は自由に設計図を書くことができるので，\keytermL{分散分析}{ぶんさんぶんせき}や\keytermL{t検定}{tけんてい}など平均値差だけにこだわっている方法をみると，そちらがいかにも窮屈そうにしていると思えるかもしれません。講義を通じて得た知識や技術で，自由にデータを調理していただければと思います。
Modeling, which allows you to freely draw design diagrams, may make you feel that methods that are only obsessed with average difference, such as analysis of variance and t-tests, are very restrictive. Using the knowledge and skills gained through the lectures, I hope you will be able to freely manipulate the data.

% \section{帰無仮説検定の代案}
\section{Alternatives to Null Hypothesis Testing}
% ベイジアンモデリングは，自由に書いた設計図から\keyterm{確率的プログラミング言語}{かくりつてきぷろぐらみんぐげんご}{stocastic programming language}をつかって答えを得る方法です。プログラミング言語が使えるようになるの大変です。これまでの一般的な統計ソフトでは，ボタン1つで答えを出してくれたのに，と思う人もいるかもしれません。しかしそれは，従来の分析方法がさまざまな仮定や前提に基づいた，マニュアル的アプローチを許す範囲のものに限定されていたからですし，「簡単，楽ちん」ということが「知らなくていい」ということだとすると，その考え方がさまざまな問題を引き起こしてきたことを知らなければなりません。これは\textbf{心理学における再現性問題}\parencite{Ikeda2016}として取り沙汰されており，心理学そのものに疑問符を投げかけるような事態を招いています。その原因の1つが，統計的技術の誤用・悪用にありました。
Bayesian modeling is a method of obtaining answers using a stochastic programming language from a freely written blueprint. It's hard to get the hang of programming languages. Some people might think, "In traditional statistical software, you could get an answer at the touch of a button". However, that's because conventional analysis methods were limited to those that allowed for manual approaches based on various assumptions and premises, and if the concept of "easy, simple" means "it's okay not to know", then we need to understand that mentality has caused various problems. It has been discussed as the "Reproducibility issues in psychology" and it has led to situations that cast doubt on psychology itself. One of the causes was the misuse and abuse of statistical techniques.

% 心理学ではながらく，帰無仮説検定によって理論を積み重ねてきました。帰無仮説検定では$p$値が5\%より小さければ良し，という「ここだけ見ておけば良い」というようなマニュアル的基準があり，差があればすごい効果が発見されたぞ！と喧伝するということで進んできた側面があります。すでに述べたように，この\keytermL{有意差}{ゆういさ}を出すために要因計画の方にこそ心理学者は注力してきたのですが，$p<0.05$になりさえすれば良いという表面的な理解しかしていないマニュアル人間ができあがってしまうことが，問題を大きくしたのかもしれません。
In psychology, theories have been built up for a long time through null hypothesis testing. In null hypothesis testing, there are manual-like criteria such as "it's okay if the p-value is less than 5%", and there has been a trend to move forward by advertising that a great effect has been discovered if there is a difference. As already mentioned, psychologists have been focusing on factor planning to produce this "significant difference", but the fact that people who understand only superficially that it is okay if p < 0.05 becomes a manual human, may have exacerbated the problem.

% ベイズ統計によるアプローチの場合，データがどのように出てきたのか，事前分布はどうなのか，といったことを考えないわけにはいきません。自由に設計図を描くことができるというのは，裏を返せば，最初は白紙でしかないということであり，きちんと動く統計ツールを作るためには細部まですべて記述しなければならないのです。帰無仮説検定では，極端な場合「データ生成には正規分布が仮定されている」ということを知らなくても，$p$値がどうなるかだけをみる事はできてしまうのに，です。このように，ベイズ統計を使う利点の1つは，前提や仮定についてすべて自覚的に，明示的に記述しなければならないというところにあり，これが悪用・誤用を生みにくくする安全装置となっている面があります。
In the case of the Bayesian statistics approach, one has to consider how the data came out and what the prior distribution is like. The freedom to draw blueprints implies that you initially start with a blank slate and must describe everything in detail to create working statistical tools. In null hypothesis significance testing, you can just look at the p-values, even in extreme cases, without knowing that the data generation assumes a normal distribution. One of the advantages of using Bayesian statistics is that it forces you to be conscious and explicit about all assumptions and prerequisites. This serves as a safety mechanism, reducing misuse and misinterpretation.

% 最近ではJASP\parencite{JASP2021}という，GUIでベイズ分析もできるソフトウェアが出てきました。帰無仮説検定と同じような分析を，ベイズでやったらどうなるか，すぐに試すことができます。この便利なアプリではStanやJAGSなどとは違って，どのような前提があるかを自覚しなくても，分析結果を得る事はできます。ですから，今挙げた「安全装置」としてのベイズ統計の利点は，簡単に外れてしまう日が来るかもしれませんし，その方がユーザにとってはいいかもしれません。
Recently, the JASP software, which can perform Bayesian analysis with a GUI, has been introduced \parencite{JASP2021}. You can quickly try what would happen if you conducted the same analysis as the null hypothesis test in Bayesian. With this convenient app, you can obtain analytical results without being aware of any underlying assumptions, unlike Stan or JAGS. Therefore, the advantages of Bayesian statistics as a “safety device” that I mentioned earlier may one day be easily overridden, and that might be better for users.
% この講義では第\ref{x19_modeling1}講から第\ref{x22_modeling4}講まで，帰無仮説検定の代わりにモデリングアプローチをする例を紹介してきました。
In this lecture, from lecture \ref{x19_modeling1} to lecture \ref{x22_modeling4}, we have introduced examples of taking a modeling approach instead of a null hypothesis testing.
% さて，ではこうした帰無仮説検定の代わりとして，ベイズ統計を使う利点はどこにあるでしょうか。
So, where would be the advantages of using Bayesian statistics as an alternative to these null hypothesis tests?

% もちろんいくつか考えられますので，箇条書き的に解説してみたいと思います。
Of course, several things come to mind, so I would like to try explaining them in a bullet-point format.

% \paragraph{幾重にも組み合わさる仮定や補正からは無縁} たとえば二群の平均値の差を検定する場合，データが正規分布に従っているかどうかを検定し，また二群の分散が等しいかどうかも検定して，問題なければt検定に進みます。もし分散が等しくないということになれば，Welchの補正をしなければならない，とマニュアルにはあります。またたとえば，WithinのANOVAをするときも，球面性の検定を行なってデータが分析の前提となる仮定にあっているかどうかを判定しなければなりません。主効果が見られた，交互作用が見られたとなれば，今度はどこに差があるのかをみるために--検定の繰り返しを避けるために，下位検定を行います。
The text you provided translates to:

\paragraph{Unrelated to the multitude of intertwined assumptions and corrections} For instance, when testing the difference in means between two groups, you must first test whether the data follows a normal distribution, and then whether the variances of the two groups are equal. If these tests are passed, you proceed with a t-test. If the variances are found to not be equal, the manual dictates that you must apply Welch's correction. Similarly, when performing a within-ANOVA, you must carry out a sphericity test to determine whether the data meets the assumptions underlying the analysis. If main effects or interactions are observed, you must then perform subsidiary tests to identify where the differences lie -- in order to avoid repeated t-tests.

% このように，仮説検定をする場合は条件にあっているかどうかを検定し，あっていなければ補正をするなどの手続きが必要です。同時に，仮説検定を同一のデータに何度も繰り返す事は，5\%水準という\keytermL{タイプ1エラー}{たいぷ1えらー}を生じる確率が適切に制御できなくなるため，本来なら実験単位ではなく「一連の分析単位」でこれを調整しなければなりません。正直なところ，ここまで厳密な調整が，あらゆる心理学研究できちんと行われているかと言われれば，その答えはNoなのですが。
Thus, when conducting hypothesis testing, it is necessary to check whether the conditions are met and, if not, correct them, among other steps. At the same time, repeatedly performing the hypothesis test on the same data can result in a failure to properly control the probability of generating a so-called 5% level, or type 1 error. Ideally, this should be adjusted not on a per-experiment basis, but on a "per-sequence of analysis" basis. To be honest, if asked whether such rigorous adjustments are being properly made in all psychological research, the answer would be no.

% ベイズ統計の場合，こうした問題は非常に軽微なものに変わります。たとえば二群の分散が等しいかどうかについては，「等しくないモデル」を考えれば良いのです。具体的にいうと，二群の分散が等しいモデルは，
In the case of Bayesian statistics, such problems become very minor. For example, for whether the variances of the two groups are equal, all we need to do is consider the "unequal model". Specifically, for a model where the variances of the two groups are equal,
% $$ X_i \sim N(\mu_1,\sigma), Y_i \sim N(\mu_2,\sigma)$$
As an AI, I can inform you that this is a mathematical expression, not a Japanese text. This expression represents that X_i and Y_i are variables that follow a Normal distribution. X_i follows the Normal distribution with mean μ1 and standard deviation σ. Similarly, Y_i follows the Normal distribution with mean μ2 and standard deviation σ.
% であり，等しくないモデルは
The model is fixed and not equal.
% $$ X_i \sim N(\mu_1,\sigma_1), Y_i \sim N(\mu_2,\sigma_2)$$
To translate the given LaTeX code, it's a mathematical expression and not Japanese text, it states that:

"X_i follows a normal distribution with mean mu_1 and standard deviation sigma_1, Y_i follows a normal distribution with mean mu_2 and standard deviation sigma_2."
% とするだけです($\sigma$に添字がある，すなわち違う大きさとして推定する)。球面性の検定についても分散共分散行列の要素それぞれを推定するだけであり，その数値に事前に定められた制約はありません。
It only needs to be assumed (a subscript is on $\sigma$, that is, it is estimated as a different size). The test for sphericity is simply an estimation of each element of the variance-covariance matrix, and there is no preset constraint on the value.

% またたとえば，タイプ1エラーの制御についても，そもそもタイプ1エラーという発想がありませんから存在しないのです。データがあって，平均構造を線形モデルで表現したのであれば，そこから出てくる事後分布は1つしかありません。その単一のパラメータ同時事後分布をどのように切り分けようとも，ある判断のエラーに確率が伴う事はないのです。
For example, regarding the control of Type 1 errors, there is no such concept as a Type 1 error in the first place. There is data, and if the mean structure is represented by a linear model, there is only one posterior distribution that comes out of it. No matter how you divide that single parameter joint posterior distribution, there is no chance of error in any judgment.

% 帰無仮説検定を厳密に行うには，下位検定を事前に計画してどこにどのような差があると考えるかも準備しておく必要があります。帰無仮説検定は勝敗を決める手法ですから，ゲームのルールは事前に決めなければならないのです。ゲームの設定如何によっては，判定結果が変わってくることがあるからです。これに対してベイズ統計的アプローチであれば，結果(事後分布)は決まっているので，下位検定を事前に計画する必要はなく，事後分布をどう切り分けようと分析者の自由です。
To conduct a null hypothesis test strictly, it is necessary to plan a subordinate test in advance and prepare where and what kind of difference is expected. Since null hypothesis testing is a method of determining wins and losses, the rules of the game must be determined in advance. This is because the judge's decision can vary depending on the setting of the game. In contrast, in the Bayesian statistical approach, the outcome (posterior distribution) is predetermined, so there is no need to plan a subordinate test in advance, and it is up to the analyst to divide the posterior distribution as he or she pleases.
% \paragraph{停止規則やサンプルサイズの問題が生じない}
\paragraph{Issues such as  stop rules and sample size do not arise}
% 同様のことは，サンプルサイズの設計についても言えます。帰無仮説検定の場合は，ゲームのルールを事前に決める必要があると言いました。このルールの中には，サンプルサイズも含まれています。サンプルサイズが大きくなればなるほど，有意であると判定される可能性は上がっていきます。帰無仮説検定における，最もやってはいけない研究実践のひとつは，データを取りながら検定を繰り返し，有意になったらデータを取るのをやめる，というものです。検定を繰り返すと，誤った判断をしてしまう可能性が増えますし，サンプルサイズが大きくなると有意になりやすくなります。つまり「勝つまでゲームをやめない」というのは卑怯な方法なのです。
The same can be said about the design of sample size. In the case of null hypothesis testing, I stressed that the rules of the game need to be set in advance. This includes sample size. The larger the sample size, the higher the probability of being deemed significant. One of the most forbidden practices in null hypothesis testing is to keep testing as you collect data, and stop collecting data when it becomes significant. Repeating tests increases the risk of making a wrong decision, and as the sample size gets larger, it becomes easier to become significant. In other words, "not quitting the game until you win" is an unfair method.

% しかし実際は，数人相手にデータをとってみたけど有意差はでなかった，でも惜しいから頑張ってデータを増やした，結果としてある日有意な差になったので，サンプルサイズをちゃんと記載して論文として提出した，ということがあるわけです。努力して真実に辿り着くのはいい話ですが，この例は努力したことが結果の捏造になっていたかもしれないという，より悲しい話になってしまいます。
But in reality, I tried collecting data from several people, but there was no significant difference. Nevertheless, out of regret, I worked hard to increase the data. As a result, one day there was a significant difference, so I properly stated the sample size and submitted it as a paper. This could happen . It's a good story to make an effort and reach the truth, but in this case, the effort may have resulted in falsification of results, which is a more sad story.

% 帰無仮説検定は事前にサンプルサイズを決めなければいけません。このサイズ，この基準で勝負するぞ，というルールがあってからの一発勝負なのです。どこでデータの収集を止めるかというのは，\keyterm{停止規則}{ていしきそく}{Stopping Rule}といって，これまた事前に決めておかねばならないことなのです。事後的に，勝負に勝ったからこのやり方でよかったんだろう，と考えるのは間違っているのです。
The null hypothesis significance testing requires determining the sample size in advance. It's like taking a one-shot game under rules such as "We will play based on this size and this standard". When and where to stop collecting data, referred to as the "Stopping Rule", must also be determined beforehand. It is a mistake to think after the fact that because you won the game, the approach was probably correct.

% これに対して，ベイズモデルの場合はこうした問題が生じません。最後が確率的な判断にならないので，こうした問題が生じないのです。データが蓄積されていくことは，徐々に確信度が高まっていくことでもあります。差があるのかないのか，という判断をするにあたって，よりその違いが明確になっていくだけなのです。
In contrast, no such problems arise with the Bayes model. These problems do not arise because the final decision is not probabilistic. The accumulation of data also increases confidence gradually. Whether there is a difference or not, the judgment only becomes clearer as the difference becomes more evident.

% \paragraph{より柔軟な仮説を立てることができる}
It becomes possible to make more flexible hypotheses.
% 帰無仮説検定は，帰無仮説と対立仮説という2つの仮説の比較判断になります。ここで帰無仮説は「差がない」「相関がない」といったものになります。たとえば二群間の母平均を比較する場合，帰無仮説は$\mu_1 = \mu_2$になります。
The null hypothesis test is a comparative judgment of two hypotheses, the null hypothesis and the alternative hypothesis. The null hypothesis here would be things like "there is no difference" or "there is no correlation". For example, when comparing the population mean between two groups, the null hypothesis would be $\mu_1 = \mu_2$.
% しかし実際には，差がないということを主張したい場合もあるかもしれません。帰無仮説に「差がある」というのをおくことができないのは，「あるというのは，どれぐらいあればいいのか」という問題がすぐに浮上するからです。たとえば母平均の差が$\mu_1 - \mu_2 = 0.1$なのか，$\mu_1 - \mu_2 = 0.11$なのか，$\mu_1 - \mu_2 = 0.111$なのか・・・言い出すとキリがないですね。$p$値は帰無仮説のもとで計算されるものですから，差があるという仮説は無限に作れてしまうので，検証すべき帰無仮説も無限に膨れ上がってしまいます。差がないという状態は$\mu_1 - \mu_2 = 0$と一意に定まるので，帰無仮説を設定することが可能なのです。
However, in reality, there may be times when you want to argue that there is no difference. The reason why we cannot set "there is a difference" as the null hypothesis is because the problem of "how much should there be to say that there is" arises immediately. For example, is the difference in the mother's average $\mu_1 - \mu_2 = 0.1$, or $\mu_1 - \mu_2 = 0.11$, or $\mu_1 - \mu_2 = 0.111$...can be said endlessly. The $p$-value is calculated under the null hypothesis, so if there is a hypothesis that there is a difference, it can be created infinitely, so the null hypothesis to be verified will also inflate infinitely. The state of no difference can be uniquely determined by $\mu_1 - \mu_2 = 0$, so it is possible to set the null hypothesis.

% では帰無仮説を積極的に支持すればよいではないか，と思われるかもしれませんが，それができません。判定結果の$p$値は帰無仮説のもとで計算していますから，「帰無仮説のもとで考えるとおかしな結果になった」という背理法が成り立たないのです。「帰無仮説のもとでおかしな結果にならなかった」というのは，帰無仮説以外の可能性を除外できるものではありません。
You might think it would be nice to actively support the null hypothesis, but you can't. As the p-value of the decision result is calculated under the null hypothesis, we cannot form an indirect proof like "It will lead to an odd result if we think under the null hypothesis." The statement "No strange results were obtained under the null hypothesis" doesn't enable us to exclude any possibilities other than the null hypothesis.

% このように，帰無仮説検定では「効果がなかった」とか「関係がなかった」とは言えず，せいぜい「効果がないとは言えない」というにとどまるのでした。
In this way, in a null hypothesis test, we cannot say "there was no effect" or "there was no relationship", at most we can only say "we cannot say that there was no effect".

% これに対して，ベイズ統計では\keytermL{モデル比較}{もでるひかく}の観点から検証することになります。帰無仮説は$\mu_1=\mu_2$というモデル，対立仮説は$\mu_1 \neq \mu_2$というモデルです。言い換えれば，帰無仮説は
In contrast, in Bayesian statistics, validation will be from the viewpoint of \keytermL{model comparison}. The null hypothesis is a model where $\mu_1=\mu_2$, and the alternative hypothesis is a model where $\mu_1 \neq \mu_2$. In other words, the null hypothesis is a model where the means are equal under all conditions.
% $$ X_i \sim N(\mu_1,\sigma_1), Y_i \sim N(\mu_1,\sigma_2)$$
As a language translator, I'm unable to translate mathematical notation or codes. But, I can tell you that the provided notation is a representation of two statistical distributions in probability theory and statistics:

1. Xi follows a normal distribution with mean mu1 and standard deviation sigma1.
2. Yi follows a normal distribution with mean mu1 and standard deviation sigma2.
% というモデルであり，対立仮説は
This model is, the alternative hypothesis is
% $$ X_i \sim N(\mu_1,\sigma_1), Y_i \sim N(\mu_2,\sigma_2)$$
This Japanese text is actually a mathematical expression known as a Normal distribution in statistics. The translation would be:

"Variable X_i follows a normal distribution with mean µ_1 and standard deviation σ_1. Variable Y_i follows a normal distribution with mean µ_2 and standard deviation σ_2."

% というモデルだとしているに過ぎないからです。
This is simply because it is considered as a model.

% これまでみてきたように，ベイズ統計のアプローチで平均値の差を考える場合は，$\mu_1 - \mu_2$のような差の分布を\keytermL{生成量}{せいせいりょう}でつくり，その大きさをそのまま吟味できます。どうしても「差があるのか，ないのか」という判断がしたければ，\keyterm{実質的に等価な範囲}{じっしつてきにとうかなはんい}{Region Of Practical. Equivalence; ROPE}($\to $ セクション\ref{ROPE}, Pp.\pageref{ROPE})を事前に設定し，その区間に差の分布が入ってくるかどうかを判断するのでした。そもそも差も分布しているわけですから，どれぐらい重複しているかという幅で考えるほかないわけです。帰無仮説検定の文脈においても，\keyterm{効果量}{こうかりょう}{Effect Size}とよばれる\textbf{標準化された差}の大きさを算出して考えますが，このように「どれぐらい違いがあれば差があると言えるか」という程度問題が前面に出てくることが重要なのです。5\%水準のように機械的に判断するのではなく，領域固有の知識(ドメイン知識)に基づいて，実質的な差の大きさをケースバイケースで判断することが重要です。もちろん機械的に手続きを進められないことは不便かもしれませんが，不便だからといって真実を曲げるようでは本末転倒です。
As we have seen so far, when considering the difference in mean values in the Bayesian statistical approach, we generate a distribution of differences such as $\mu_1 - \mu_2$, and scrutinize its magnitude. If you really want to make a judgment about whether there is a difference or not, you should determine a Region Of Practical Equivalence (ROPE) beforehand and decide whether the difference distribution falls within this range (→ Section \ref{ROPE}, Pp.\pageref{ROPE}). After all, since differences also distribute, we have no choice but to consider the extent of overlap. In the context of Null hypothesis testing, we calculate and consider the magnitude of the standardized difference called Effect Size, but it is crucial to consider how much of a difference can be said to be a difference. Instead of making a mechanical judgment at a 5% level, it is important to determine the size of the substantial differences case by case, based on domain-specific knowledge. Although it may be inconvenient that you cannot proceed mechanically, it would be putting the cart before the horse to distort the truth just because it's inconvenient.


% また生成量を使ったアプローチのところで解説したように，パラメータについて考えるだけでなく，そのモデルから生成されるであろうデータについて仮説を立てたり検証したりできるのも，ベイズ統計的アプローチの利点です。実際どの程度の差があると意味があるのかについて，具体的な数値シミュレーションで考えることができるからです。パラメータが$\delta$以上の差があれば良いとか，二群のデータが$D$以上違ってくる確率はどれぐらいか，といった検証の仕方は，仮想空間上の$p$値よりも具体的で意味のある情報をもっています。パラメータやデータについて，「同じかどうか」以上の情報を検証できることは，ベイズ統計の長所といえるでしょう。
As explained in the section on using generative quantities, another advantage of the Bayesian statistical approach is that it allows us to hypothesize and test not just about the parameters, but also about the data that would be generated from the model. This is because we can consider through concrete numerical simulations how much difference would be meaningful. Testing methods such as whether a difference of $\delta$ or more in the parameters would be good, or what is the probability that the data from two groups differ by $D$ or more, carry concrete and meaningful information compared to p-values ​​in virtual space. Being able to test more than just "whether they are the same or not" about parameters and data can be said to be an advantage of Bayesian statistics.

% \section{モデル比較}
\section{Model Comparison}
% 先ほど，\textbf{モデル比較}ということに言及しました。このことについて，最後に少し触れておきたいと思います。
I mentioned something about \textbf{model comparison} a little while ago. I would like to touch on this a little bit at the end.
% ベイジアンモデリング，すなわちモデルをたててベイズ統計学的に推定するなかで，そのモデルの優劣を考える方法があります。
There is a method of evaluating the superiority or inferiority of a model in Bayesian modeling, that is, in estimating it statistically based on Bayes' theorem.

% もう一度ベイズの公式を見てみましょう。ベイズの公式は次のようなものでした。
Let's take another look at Bayes' theorem. Bayes' theorem was as follows.

\[ p(\theta| D) = \frac{p(D| \theta)p(\theta)}{p(D)} \]

ここで右辺の分母は$p(D)$，つまりデータの確率をあらわすもので\keyterm{周辺尤度}{しゅうへんゆうど}{marginal likelihood}と呼ばれます。
これはどのように計算されるのでしょうか。尤度は$p(D| \theta)$で表されますが，$P(D)$はこのパラメータが無くなったものになります。どのようにしてパラメータをなくすかというと，パラメータのありそうな可能性をすべて足し合わせることで\keytermL{周辺化消去}{しゅうへんかしょうきょ}すれば良い，ということになります。

たとえば性別と学年のクロス表があったとして，ランダムに選んだ1人の学生が女性である確率は，すべての学年の女性の確率を足し合わせたものになるわけです。記号で考えると，条件付き確率$P(D| \theta)$において，$\theta$の可能性を全部考えれば$P(D)$だけになるということです。これはパラメータ$\theta$が離散的な例ですが，連続的な場合は
$$ p(D) = \int p(D | \theta ) d\theta $$
のように積分で考えればよいでしょう。

これはモデルに含まれるあらゆるパラメータのパターンを網羅して，データが得られる確率を計算したことになりますから，モデルが表現できる世界の中でどれぐらいデータを捕まえることができたか，ということを表しているとも言えます。
ここでモデルが複数ある場合を考えましょう。あるモデルを$\mathcal{M}$と表現することにして，第一のモデルを$\mathcal{M}_1$,第二のモデルを$\mathcal{M}_2$，と表したとします。
すると事後分布も「あるモデルを仮定した時の事後分布」ということになりますから，次のようにすべて「モデルのもとでの」という条件付き確率で表現できます。

$$ p(\theta | D,\mathcal{M}_i) = \frac{p(D | \theta,\mathcal{M}_i)p(\theta | \mathcal{M}_i)}{p(D| \mathcal{M}_i)} $$

ここで右辺の分母を見ると，$p(D| \mathcal{M}_i)$，つまり「モデル$i$のもとでデータが得られる確率」を表していることになります。もしこの確率が高ければ，このモデルはこのデータを生む確率が高いわけです。逆にもしこの確率が低ければ，このモデルはこのデータを生む確率が低いということになります。この周辺尤度は，モデルがデータを支持する度合いであり，証拠(\keytermL{エビデンス}{えびでんす})の強さとも言えるわけです\footnote{これを考えると，どんなデータでも作れる幅広いモデルを考えておけばいいじゃないか，と思うかもしれませんが，モデルを広く考えすぎるとそこから出てくるデータはそのごく一部でしかなく，かえってエビデンスとしては弱くなってしまいます。なんでも予測できるモデルは何も予測していないことと同じなのです。}。

その上で，次のような数字を考えます。

\[ BF_{12} = \frac{p(D| \mathcal{M}_1)}{p(D| \mathcal{M}_2)} \]

% この数字は\keyterm{ベイズファクター}{べいずふぁくたー}{Bayes Factor}と呼ばれるもので，モデル2に比べモデル1がどれほどデータに支持されているかを示す数字です。この数字が1.0より大きければ，モデル1はモデル2に比べて相対的に強くデータに支持されているということになります。逆に1.0を下回ると，モデル2のほうがモデル1よりもいいね，ということです。モデル同士の比による表現ですので，$BF_{21}$としても構いません。ここでモデル1が対立仮説，モデル2が帰無仮説だとすると，どちらがよりデータをうまく表しているか，どちらがよりデータに支持されているかがわかるわけです。最近では，帰無仮説検定の文脈でもBFを併記して，どちらのモデルが良いかを示すこともありますし，この方法ですと帰無仮説(というか差がないというモデル)を積極的に支持する，ということもできます。
This number is called the \keyterm{Bayes Factor}, and it shows how much model 1 is supported by the data compared to model 2. If this number is greater than 1.0, it means that model 1 is relatively strongly supported by the data compared to model 2. On the other hand, if it falls below 1.0, it means model 2 is better than model 1. Because it's a comparative representation between models, you can also notate it as $BF_{21}$. Assuming here that model 1 is the alternative hypothesis and model 2 is the null hypothesis, we can see which model is better represented by the data and which is more supported by the data. Recently, it has also become common to annotate BF in the context of null hypothesis testing to show which model is better, and with this method, we can actively support the null hypothesis (or rather, the model suggesting there is no difference).

% この手法はベイズ推定するものであれば一般的に通用する考え方ですので，1つのデータに対していろいろなモデルを考えついたとしても，ベイズファクターを計算して相互に比較できるわけです。すなわち，ベイズのモデル比較はベイズファクターという共通の基盤を手にしたとも言えます。
This is a concept that is generally applicable if you are doing Bayesian estimation, and even if you come up with various models for a single data set, you can calculate the Bayes factor and compare them with each other. In other words, you could say that Bayesian model comparison has a common foundation in the Bayes Factor.

% そんな優れた方法があるのであれば，なぜもっと早く教えてくれなかったのか，という声が聞こえてきそうです。実は理屈上ではこの通りなのですが，この方法には実践上の問題が残されているのです。それは，\underline{周辺尤度の計算が難しい}というものです。すべてのパラメータについて，考えられるすべての可能性を計算して足し上げるわけですから，それはそれは大変な計算になります。有効な近似計算の方法として，全パラメータの全確率空間を計算し尽くすかわりに，確率が高そうなところを効率よくサンプリングして推定値を得る\keyterm{ブリッジ・サンプリング}{ブリッジ・サンプリング}{Bridge Sampling}という手法が考えられており，その計算をするRのパッケージ\parencite{bridgesampling}が提供されていますから，今後こうしたモデル比較も増えてくるでしょう。
It might remind you the question: "Why didn't they tell me such a great method much sooner?" In fact, in theory, it is indeed as stated, but there are practical issues left with this method. That is, \underline{the calculation of marginal likelihood is difficult}. After all, it means adding up all possible values for all parameters, which is a very difficult calculation. An efficient approximation method called Bridge Sampling is considered as an effective approximate calculation, which obtains estimates by efficiently sampling places where the probability is likely to be high, instead of calculating the entire probability space for all parameters. Since the R package, which performs such calculations, is provided, more model comparisons will be expected in the future.

% またモデル評価の方法として，\keytermL{情報量規準}{じょうほうりょうきじゅん}によるアプローチも考えられています。これは「良いモデルとは良い予測をするモデルである」と考えるところから始まっており，平均的に良い予測をするモデル程度を指標化するものです。ベイズの定理を基本としながら，データを真のモデルから生成されるサンプルと考えたり，確率分布同士の近さを\keyterm{カルバック・ライブラー情報量}{かるばっく・らいぶらーじょうほうりょう}{Kullback-Leibler divergence}で表現するなど，非常に高度な数学的体系が確立されています。詳しくは\textcite{Hamada201912}など丁寧な解説書があり，Rパッケージなどで簡単に計算できるようになっています\parencite{loo}。数学的にやや複雑な内容を含んでいますので，本講義の中では扱いきれんませんが，より進んだ研究や利用のために言葉の紹介だけはしておきます。
As a method of model evaluation, an approach based on the \keytermL{information criterion}{じょうほうりょうきじゅん} is also considered. This starts from the idea that "a good model is one that makes good predictions", and quantifies the degree to which a model makes good predictions on average. While based on Bayes' theorem, it involves a highly advanced mathematical system that considers data as samples generated from the true model and represents the closeness between probability distributions in terms of the \keyterm{Kullback-Leibler divergence}{かるばっく・らいぶらーじょうほうりょう}{Kullback-Leibler divergence}. Detailed explanations can be found in \textcite{Hamada201912} and others, and calculations can be easily done using R packages \parencite{loo}. It includes somewhat complex mathematical content, so we can't cover it all in this lecture, but I'll at least introduce the terms for further research and use.


% \section{おわりに}
\section{In Conclusion}
% 本講ではここまで，ベイズ統計と従来の統計的分析を比較して議論してきました。ベイズ統計の方がメリットが大きいような解説をしてきましたが，これは筆者の好みもかなり含まれるものです。帰無仮説検定はなにも悪いものではなく，前提や仮定をしっかり守れば，一般的なソフトウェアで誰にでも結論を出すことができるという点で，とても良いものなのです。前提や仮定が多くて面倒だから，マニュアル人間ができてしまうというのは，あくまでもユーザ側の問題であって方法論の問題ではありません。ベイズ統計だって，モデルが間違っていれば結論は間違ったものになりますし，なによりデータに合わせてモデルを逐一設計図から書き起こすのは，大変コストがかかることでもあるからです。本来，これら2つの流儀は相反するものではなく，それぞれ依って立つ理論や考え方があって，必要に応じて選べるようになるのが一番です。
In this lecture, we have been comparing and discussing Bayesian statistics and conventional statistical analysis. Although we have been explaining that Bayesian statistics has greater merits, this includes quite a bit of the author's preference. The null hypothesis test is not a bad thing at all, and if the assumptions and prerequisites are properly respected, anyone can reach a conclusion using general software, which is a very good thing. If there are many assumptions and prerequisites and it's troublesome, and people just follow the manual, this is a problem for the users, not a problem with the methodology. Even with Bayesian statistics, if the model is wrong, the conclusion will be wrong, and designing a model from scratch for each piece of data is a high-cost task. Ideally, these two methods should not be mutually exclusive, each has its own theory and perspective, and it's best to be able to choose as needed.

% またベイズ流の研究アプローチ，とくにモデル比較については，まだまだ議論が活発に行われている段階であり，とくに心理学研究にこうしたアプローチがどれほど有用かについては，いまだに結論が出ていないところでもあります。
The Bayesian approach to research, especially model comparison, is still under active discussion, and there is no conclusion yet on how useful this approach is for psychological research.
% パラメータの推定値で考えるのが良いのか，モデル比較をして議論を進めるのかについても，まだまだ定番の方法というのは見つかっていません。
Whether it's better to consider using parameter estimates or advance discussions by comparing models, a standard method for these is yet to be discovered.

% そもそも心理学の知見を数式に落として\keytermL{モデリング}{もでりんぐ}できるようにするためには，心理学そのものの理論的発展や，人間についてのより深い理解も必要です。
In order to model psychological findings into mathematical formulas, it is necessary to develop psychological theories themselves and to gain a deeper understanding of humans.
% 人間の行動と予測が完全にできるようなモデルや一貫した理論体系は，まだまだ遥か先の目標に過ぎないというのが現状でしょう。
The current situation is that a model or consistent theoretical system that can fully predict human behavior is still a far-off goal.
% それでも有用なツールを武器に，少しずつ知見を積み重ねていくことができれば，心理学をより楽しく学べるのではないかとおもいます。Enjoy!
"Even so, if you can accumulate knowledge little by little with useful tools as weapons, I think you will be able to enjoy studying psychology even more." Enjoy!



\end{document}
